\MakeAbstract{
    随着智能建造与绿色低碳理念的发展，传统混凝土施工方式的局限性逐渐显现，挤出式3D打印混凝土（3DCP）技术为建筑工业化与绿色化转型提供了新的路径。然而，现有研究对喷嘴内部流场及整形板几何参数影响机制的系统分析仍较有限。本研究建立了含完整整形板结构的挤出式3D打印混凝土喷嘴三维计算流体动力学（CFD）模型，采用Herschel-Bulkley本构关系描述混凝土浆体的非牛顿黏塑性流变行为，对20个参数化工况进行了系统数值实验。研究以整形板倾斜角（$\alpha$）、延伸长度（TL）、喷嘴收缩角（$\theta$）及竖向高度（L）为设计变量，以基底接触压力、喷嘴壁面剪切应力、压力传递效率、压力非均匀性指数与入口压力为评价目标，系统揭示了各几何参数对挤出性能的影响规律。

    研究结果表明，整形板这一几何结构能够显著改善打印混凝土的挤出与压实性能，其中倾斜角对基底接触压力影响最为显著，板长次之，喷嘴倾斜角度具有协同正向作用，而竖向高度的统计影响不明显。在此基础上，研究进一步推导了喷嘴离地高度、整形板几何参数与混凝土层高之间的解析关系，并将其嵌入五目标Pareto-NSGA-II多目标优化框架，对离散Pareto最优解与代理模型最优解进行了比较分析。结果发现，入口压力与基底压力之间存在明显共线性，反映了约束挤出系统中能耗与压实效果之间的内在权衡关系。研究结果可为3DCP喷嘴与整形板的参数优化及工程应用提供理论依据与设计参考。
}{3D打印混凝土；计算流体动力学；Herschel-Bulkley模型；整形板；多目标优化；Pareto分析}

\MakeAbstractEng{
    Limitations of traditional concrete construction methods have become increasingly apparent under the development of intelligent construction and the newly-rising low-carbon engineering concepts. Extrusion-based 3D concrete printing (3DCP) technology provides a novel solution for the industrialization and sustainable transformation of the construction industry. However, existing studies on the internal flow field of nozzles and the influence mechanisms of trowel geometric parameters remain limited. This study presents a three-dimensional computational fluid dynamics (CFD) model of an extrusion-based 3D concrete printing (3DCP) nozzle integrated with an inclined trowel structure. Herschel-Bulkley constitutive model is incorporated to describe the Non-Newtonian viscoplastic rheological behaviour of fresh concrete. A total of 20 CFD simulation cases were analysed to investigate the effects of trowel inclination angle ($\alpha$), trowel extension length (TL), nozzle contraction angle ($\theta$), and nozzle vertical height (L) on extrusion performance. Substrate contact pressure, nozzle wall shear stress, pressure transmission efficiency, pressure non-uniformity index, and inlet pressure were selected as evaluation metrics.

    The results indicate significant improvements in both compaction and extrusion performance with the trowel structure. Among all parameters, the trowel inclination angle exhibits the most significant influence on substrate contact pressure, followed by the trowel length, while the nozzle contraction angle shows a positive synergistic effect. In contrast, the statistical influence of nozzle vertical height is not significant. A five-objective Pareto-NSGA-II optimization framework was developed by integrating analytical relationships among nozzle standoff height, trowel geometric parameters, and printed layer height. Comparative analyses between discrete Pareto-optimal solutions and surrogate-model-based optimal solutions were conducted. The results also reveal a strong collinearity between inlet pressure and substrate pressure, reflecting the intrinsic trade-off between energy consumption and compaction performance in constrained extrusion systems. The findings of this study provide theoretical support and engineering guidance for the parametric optimization and practical design of 3DCP nozzles and trowel structures.
}{3D concrete printing; computational fluid dynamics; Herschel--Bulkley model; trowel; multi-objective optimization; Pareto Front analysis}
